\begin{table}[htbp]
  \centering
  \caption{Input parameter bounds for AASM Benchmark Case 3, adapted from \citet{bekemeyer2025}.}
  \label{tab:param_bounds}
  \begin{tabular}{llcc}
    \hline
    \textbf{Parameter} & \textbf{Symbol} & \textbf{Lower Bound} & \textbf{Upper Bound} \\
    \hline
    Mach number & $M$ & 0.2 & 0.7 \\
    Angle of attack & $\alpha$ (deg) & $-3.0$ & 5.0 \\
    Reynolds number & $Re$ & $1\times10^6$ & $6.5\times10^6$ \\
    CST1 (shared) & $X_{u,1} = X_{l,1}$ & 0.0644 & 0.1932 \\
    CST2 (upper) & $X_{u,2}$ & 0.0688 & 0.2064 \\
    CST3 (upper) & $X_{u,3}$ & 0.0961 & 0.2883 \\
    CST4 (upper) & $X_{u,4}$ & 0.0961 & 0.2882 \\
    CST5 (upper) & $X_{u,5}$ & 0.1010 & 0.3030 \\
    CST6 (lower) & $X_{l,2}$ & 0.0680 & 0.2039 \\
    CST7 (lower) & $X_{l,3}$ & 0.1126 & 0.3377 \\
    CST8 (lower) & $X_{l,4}$ & 0.0381 & 0.1143 \\
    CST9 (lower) & $X_{l,5}$ & $-0.0586$ & $-0.0195$ \\
    \hline
  \end{tabular}
  \begin{tablenotes}
    \small
    \item CST1 is shared between upper and lower surfaces.
    \item Designs are derived from the RAE 2822 baseline using CST perturbation.
    \item 597 total samples (497 train, 100 test) generated using a Sobol sequence.
  \end{tablenotes}
\end{table}
